\section{Constant and variables}
\label{sec:types:constants_variables}

Every symbol a program declares has a name, or \textit{identifier}. Identifiers
are case-sensitive -- \token{count} and \token{Count} are two different names --
and they are built from letters, digits and underscores, subject to two rules: an
identifier contains at least one letter, and no digit may appear before the first
letter. Any number of leading underscores is allowed.

So \token{__id1__} is a valid identifier, while \token{__1id__} is not: the digit
arrives before any letter does.

\notebox{
  The formal grammar of an identifier is the following: \token{(\_)* [a-zA-Z]([a-zA-Z]|[0-9]|\_)*}
}

A handful of names obey those rules and are still unavailable: the language has
reserved them for statements and operations of its own. They are the
\textit{keywords}, and there are no others than these.

\token{alias}, \token{assert}, \token{atomic}, \token{break}, \token{cast},
\token{catch}, \token{class}, \token{const}, \token{copy}, \token{cte},
\token{dcopy}, \token{def}, \token{dg}, \token{dmut}, \token{do}, \token{else},
\token{enum}, \token{expand}, \token{extern}, \token{false}, \token{for},
\token{fn}, \token{future}, \token{if}, \token{impl}, \token{in}, \token{is},
\token{lazy}, \token{let}, \token{loop}, \token{macro}, \token{match},
\token{mod}, \token{move}, \token{mut}, \token{none}, \token{null}, \token{of},
\token{panic}, \token{\_\_pragma}, \token{record}, \token{union}, \token{ref},
\token{return}, \token{spawn}, \token{static}, \token{entity}, \token{throw},
\token{throws}, \token{trait}, \token{true}, \token{typeof}, \token{\_\_test},
\token{unsafe}, \token{use}, \token{\_\_version}, \token{while}, \token{with}.


Each of them is introduced where the book needs it; none is used before it has
been explained.
